% !PRINTED 568
\weberpage{568}

\begin{center}
§ 152.\qquad \textbf{The Poles of the Function} $\wp(\mu u)$.
\end{center}

In order to establish the degree of the functions $R$ and $P$, we must count the zeros of $P$, that is, the points of infinity of $\wp(\mu u)$ which are not at the same time points of infinity of $\wp(u)$. But $\wp(\mu u):\wp(u)$ becomes infinite then and only then, when
\weq{1}{u = \dfrac{h_1\omega_1 + h_2\omega_2}{\mu}}
holds, where $h_1$, $h_2$ are integers and where $h_1\omega_1 + h_2\omega_2$, but not $(h_1\omega_1+h_2\omega_2):\mu$, is a period. If we therefore set
\weq{2}{\wp(\mu u) = \dfrac{R_\mu[\wp(u)]}{P_\mu[\wp(u)]},}
and denote by $g$ the roots of $P_\mu(x)$, then
\weq{3}{g = \wp\left(\dfrac{h_1\omega_1+h_2\omega_2}{\mu}\right).}
The number $g$ does not change when $h_1$ and $h_2$ are altered by multiples of $m=\mu\mu'$, because by this the argument of the $\wp$-functions is, according to (11), § 151, altered by a period.

Since the functions $\wp(u)$ and $\wp(\mu u)$ are not altered when we subject the period $\omega_1$, $\omega_2$ to a linear transformation with determinant 1, we do not restrict the generality by assuming that $A$ is relatively prime to $m$; and on account of the periodicity of $\wp(u)$ we may therefore also set:
\[ g = \wp\left(\dfrac{h_1\omega_1+h_2 A\omega_2}{\mu}\right). \]
On account of the homogeneity of $\wp(u)$ we may now set $\omega_1=1$, $\omega_2=\omega$, and obtain, according to § 151, (4):
\[ A\omega = \dfrac{-B+\sqrt{\Delta}}{2}, \]
and accordingly we have
\weq{4}{h_1\omega_1 + h_2 A\omega_2 = h_1 + h_2\left(\dfrac{-B+\sqrt{\Delta}}{2}\right).}
If we therefore set
\[ h_1 + h_2\dfrac{-B+\sqrt{\Delta}}{2} = \nu \]
then from (3) there results:
\weq{5}{g = \wp\left(\dfrac{\nu}{\mu}\right),}
%CONT>>
